% ============================================================
% 纯答案册·课时02-倾斜角与斜率 —— 工具/答案抽册器.py 抽册生成件（勿手改；第三档＝纯答案单独成册）
% 源件（只读）：C:/提示词/工作区/M3-第2章量产0913/成卷/导学件/课时02-倾斜角与斜率/main.tex（md5 0706a6dcbde565479d1dc10d276dac00）
%% 口径：题面不重印；逐题＝题号(印面号同正册)＋[答案]＋[详解]，值/详解照源件逐字
%       （钉值门硬断言：册值≡正册件内值≡值台账值tex，逐字全等）。册序＝正册装配序＝印面号 1..25 升序（同序同号）；键序断言基准＝题面库冻结 manifest canonical 键序。
%% 三档导出：整体 main-true／纯题 main-false（在产）｜本册＝纯答案册（本件）
%% 双档：ansbook-true.tex＝详解印本；ansbook-false.tex＝纯值速查本（详解不印）
% 编译：xelatex <壳> ×2，cwd＝本目录；sty＝qp-m3.sty 本地挂载（md5 7c3930362be8a0a2bdf21bbf8ac16573，锁校验）
% ============================================================
\documentclass[fontset=none]{ctexart}
\usepackage{qp-m3}
% —— 印面逐字符号路由补钉（承源件；− 走 TNR、▱ NSC 子块；禁改 sty 保 md5） ——
\xeCJKDeclareCharClass{Default}{"2212}%
\setCJKfamilyfont{nscblk}[Path=C:/提示词/工作区/字替对照-0909/variantF/fonts/,BoldFont=NSC-w600.ttf]{NSC-w500.ttf}%
\xeCJKDeclareSubCJKBlock{nscblk}{"25B1}%
\newcommand{\pxparallelogram}{{\CJKfamily{nscblk}▱}}%
% —— 断行微调（承母版实测档，三零门承重；\emergencystretch 不另设） ——
\xeCJKsetup{CJKglue={\hskip 0.10em plus 0.08em minus 0.05em}}%
% —— 纯答案册全线括线（长详解可跨栏断，承练习件全线括线口径；灰底盒不可跨栏断） ——
\ansblockgrayfalse
% —— 双档开关（false 档＝纯值速查：答案印、详解不印；件面开关，不动 sty 本体） ——
\newif\ifansbookdetail
\ifdefined\ansbookdetail
  \ifnum\ansbookdetail=0 \ansbookdetailfalse\else\ansbookdetailtrue\fi
\else
  \ansbookdetailtrue
\fi
% —— 页脚件名（承源件章名；件名＝<件型>答案册） ——
\renewcommand{\qpzhangming}{第二章\quad 平面解析几何}%
\renewcommand{\qpjianming}{导学件答案册}%

\begin{document}
\zhangtitle{第二章\quad 平面解析几何}
\jietitle{2.2.1 直线的倾斜角与斜率（第1课时）}
\par\glueguard{1}\addvspace{2pt}\noindent\biaoqian{【纯答案册】}{\fangsong 题面不重印\quad 印面号与正册同号同序}\par\addvspace{2pt}

\begin{multicols}{2}
\emergencystretch=1em

\begin{ansblock}[2章-练-课时02-E1]
% ans:2章-练-课时02-E1
\ansitem{1}{C}
\ifansbookdetail
\ansnote{详解}{\(O\)、\(O_3\) 都为五角星的中心点，故 \(OO_3\) 平分第三颗小星的一个角；又五角星的内角为 \(36^{\circ}\)，可知\(\angle BAO_3=18^{\circ}\)．过 \(O_3\) 作 \(x\) 轴的平行线 \(O_3E\)，则\(\angle OO_3E=\alpha\approx 16^{\circ}\)，所以直线 \(AB\) 的倾斜角约为 \(18^{\circ}-16^{\circ}=2^{\circ}\)，选C．}
\fi
\end{ansblock}

\begin{ansblock}[2章-练-课时02-E9]
% ans:2章-练-课时02-E9
\ansitem{2}{B}
\ifansbookdetail
\ansnote{详解}{\(\sin 210^{\circ}=-1/2\)，直线方程即\(-x-y-2=0\)，即\(y=-x-2\)，斜率\(k=2\sin 210^{\circ}=-1\)，故倾斜角为 \(135^{\circ}\)，选B．}
\fi
\end{ansblock}

\begin{ansblock}[2章-练-课时02-E15]
% ans:2章-练-课时02-E15
\ansitem{3}{120°}
\ifansbookdetail
\ansnote{详解}{\(l_2\) 与 \(l_1\) 垂直，则 \(l_2\) 的向上方向与 \(x\) 轴正向所成的角为\(30^{\circ}+90^{\circ}=120^{\circ}\)（在 \(l_2\) 向上方向取与 \(l_1\) 垂直且指向左上的一侧，另一直线方向与之共线）．由倾斜角的定义及\(0^{\circ}\leqslant\alpha<180^{\circ}\)，得 \(l_2\) 的倾斜角为\(120^{\circ}\)．（说明：两直线垂直时倾斜角关系的斜率判据在§2.2.3学习，本题只需用平面几何角的运算作答．）}
\fi
\end{ansblock}

\begin{ansblock}[2章-练-课时02-E16]
% ans:2章-练-课时02-E16
\ansitem{4}{(1) 60°；(2) 150°；(3) 90°；(4) 0°．}
\ifansbookdetail
\ansnote{详解}{（1）\(k=\sqrt{3}\)，\(\tan\alpha=\sqrt{3}\)，\(\alpha=60^{\circ}\)；（2）化为\(y=-(\sqrt{3}/3)x-\sqrt{3}/3\)，\(k=-\sqrt{3}/3\)，\(\tan\alpha=-\sqrt{3}/3\)，\(\alpha=150^{\circ}\)；（3）直线垂直于 \(x\) 轴，斜率不存在，倾斜角\(\alpha=90^{\circ}\)；（4）直线与 \(x\) 轴平行，倾斜角\(\alpha=0^{\circ}\)．}
\fi
\end{ansblock}

\begin{ansblock}[2章-练-课时02-E7]
% ans:2章-练-课时02-E7
\ansitem{5}{C}
\ifansbookdetail
\ansnote{详解}{三点共线\(\iff k_{AB}=k_{AC}\)．\(k_{AB}=(2-3)/(3+2)=-1/5\)；\(k_{AC}=(m-3)/(1/2+2)=(m-3)/(5/2)=2(m-3)/5\)．令\(2(m-3)/5=-1/5\)，得\(2(m-3)=-1\)，\(m=5/2\)，选C．}
\fi
\end{ansblock}

\begin{ansblock}[2章-练-课时02-E14]
% ans:2章-练-课时02-E14
\ansitem{6}{a＝1}
\ifansbookdetail
\ansnote{详解}{由斜率公式\(k=[(a+1)-(-1)]/(2-a)=(a+2)/(2-a)=3\)（须\(a\neq 2\)），去分母得\(a+2=3(2-a)=6-3a\)，\(4a=4\)，\(a=1\)，且\(a=1\neq 2\)满足条件．所以\(a=1\)．}
\fi
\end{ansblock}

\begin{ansblock}[2章-练-课时02-E10]
% ans:2章-练-课时02-E10
\ansitem{7}{(−∞,1)∪(1,+∞)}
\ifansbookdetail
\ansnote{详解}{三点能构成三角形\(\iff\)三点不共线．由\(k_{AB}\neq k_{AC}\)，即\((k-1)/(-2-3)\neq(1-1)/(8-3)\)，得\((k-1)/(-5)\neq 0\)，\(k-1\neq 0\)，\(k\neq 1\)．故\(k\)的取值范围为\((-\infty,1)\cup(1,+\infty)\)．}
\fi
\end{ansblock}

\begin{ansblock}[2章-练-课时02-E6]
% ans:2章-练-课时02-E6
\ansitem{8}{BCD}
\ifansbookdetail
\ansnote{详解}{倾斜角为钝角\(\iff\)斜率存在且为负，即\(A\)、\(B\)横坐标不等且\(k_{AB}<0\)．\(k_{AB}=(1+a-a)/[(1-a)-3]=1/(-2-a)<0\)，得\(a+2>0\)，即\(a>-2\)．又\(a=-2\)时\(A(3,-1)\)、\(B(3,-2)\)横坐标相同、斜率不存在，倾斜角\(90^{\circ}\)，不在"钝角"取值内且\(a>-2\)已排除．故\(a>-2\)，选项中\(0\)、\(1\)、\(2\)满足，选BCD．}
\fi
\end{ansblock}

\begin{ansblock}[2章-练-课时02-E2]
% ans:2章-练-课时02-E2
\ansitem{9}{A}
\ifansbookdetail
\ansnote{详解}{直线\(PA\)的斜率为\(k_{PA}=(-2+1)/(1-0)=-1\)，直线\(PB\)的斜率为\(k_{PB}=(1+1)/(2-0)=1\)．数形结合：当过\(P\)的直线\(l\)绕\(P\)从\(PA\)旋转到\(PB\)（保持与线段\(AB\)有公共点）时，其斜率介于\(k_{PA}\)与\(k_{PB}\)之间，故\(k\in[-1,1]\)，选A．}
\fi
\end{ansblock}

\begin{ansblock}[2章-练-课时02-E8]
% ans:2章-练-课时02-E8
\ansitem{10}{B}
\ifansbookdetail
\ansnote{详解}{\(\tan\alpha=k_{OA}=\sqrt{2}\)．\(m'\)的倾斜角为\(2\alpha\)，故\(k=\tan 2\alpha=2\tan\alpha/(1-\tan^{2}\alpha)=2\sqrt{2}/(1-2)=-2\sqrt{2}\)，选B．（合理性：\(\tan\alpha=\sqrt{2}>1\)，\(\alpha\in(\pi/4,\pi/2)\)，\(2\alpha\in(\pi/2,\pi)\)确为钝角，\(\tan 2\alpha<0\)．）}
\fi
\end{ansblock}

\begin{ansblock}[2章-练-课时02-E3]
% ans:2章-练-课时02-E3
\ansitem{11}{C}
\ifansbookdetail
\ansnote{详解}{\(\alpha\in[0^{\circ},180^{\circ})\)且\(\sin\alpha=3/5>0\)，故\(\alpha\)为锐角或钝角．当\(\alpha\)为锐角时\(\cos\alpha=4/5\)，\(k=\tan\alpha=3/4\)；当\(\alpha\)为钝角时\(\cos\alpha=-4/5\)，\(k=\tan\alpha=-3/4\)．故直线\(l\)的斜率是\(3/4\)或\(-3/4\)，选C．}
\fi
\end{ansblock}

\begin{ansblock}[2章-练-课时02-E4]
% ans:2章-练-课时02-E4
\ansitem{12}{B}
\ifansbookdetail
\ansnote{详解}{\(\alpha\in[0,\pi)\)，\(-1\leqslant k\leqslant\sqrt{3}\)即\(-1\leqslant\tan\alpha\leqslant\sqrt{3}\)．当\(0\leqslant\tan\alpha\leqslant\sqrt{3}\)时\(\alpha\in[0,\pi/3]\)；当\(-1\leqslant\tan\alpha<0\)时\(\alpha\in[3\pi/4,\pi)\)．所以\(\alpha\in[0,\pi/3]\cup[3\pi/4,\pi)\)，选B．}
\fi
\end{ansblock}

\begin{ansblock}[2章-练-课时02-E5]
% ans:2章-练-课时02-E5
\ansitem{13}{A}
\ifansbookdetail
\ansnote{详解}{\(k_l=(t^{2}-1)/(3-2)=t^{2}-1\)．由\(-\sqrt{2}\leqslant t\leqslant\sqrt{2}\)得\(t^{2}\in[0,2]\)，故\(k\in[-1,1]\)，即\(\tan\theta\in[-1,1]\)（\(\theta\)为倾斜角），得\(\theta\in[0,\pi/4]\cup[3\pi/4,\pi)\)，选A．（注：本选项组为本卷补写——源题为无选项填空，为入选择槽按原题数据配四选项，答案为\(A\)不变．）}
\fi
\end{ansblock}

\begin{ansblock}[2章-练-课时02-E11]
% ans:2章-练-课时02-E11
\ansitem{14}{[√3/3, √3)}
\ifansbookdetail
\ansnote{详解}{正切函数\(k=\tan\alpha\)在\([\pi/6,\pi/3)\)上单调递增，故\(\alpha\in[\pi/6,\pi/3)\)时\(\tan\alpha\in[\tan(\pi/6),\tan(\pi/3))=[\sqrt{3}/3,\sqrt{3})\)，即斜率\(k\in[\sqrt{3}/3,\sqrt{3})\)．}
\fi
\end{ansblock}

\begin{ansblock}[2章-练-课时02-E12]
% ans:2章-练-课时02-E12
\ansitem{15}{[1,+∞)∪(−∞,−√3]}
\ifansbookdetail
\ansnote{详解}{\(\theta\in[\pi/4,\pi/2)\)时，\(k=\tan\theta\in[\tan(\pi/4),+\infty)\)\allowbreak\(=[1,+\infty)\)\par\noindent \(\theta\in(\pi/2,2\pi/3]\)时，\(k=\tan\theta\in(-\infty,\tan(2\pi/3)]=(-\infty,-\sqrt{3}]\)\par\noindent 故\(k\)的取值范围为\([1,+\infty)\cup(-\infty,-\sqrt{3}]\)．}
\fi
\end{ansblock}

\begin{ansblock}[2章-练-课时02-E13]
% ans:2章-练-课时02-E13
\ansitem{16}{−2−√3}
\ifansbookdetail
\ansnote{详解}{\(l_1\)的斜率为\(-1\)，其倾斜角为\(135^{\circ}\)，故\(l_2\)的倾斜角为\(135^{\circ}-30^{\circ}=105^{\circ}\)．\(k_2=\tan 105^{\circ}=\tan(60^{\circ}+45^{\circ})=(\tan 60^{\circ}+\tan 45^{\circ})/(1-\tan 60^{\circ}\cdot\tan 45^{\circ})=(\sqrt{3}+1)/(1-\sqrt{3})=(\sqrt{3}+1)^{2}/[(1-\sqrt{3})(1+\sqrt{3})]=(4+2\sqrt{3})/(-2)=-2-\sqrt{3}\)．所以直线\(l_2\)的斜率为\(-2-\sqrt{3}\)．}
\fi
\end{ansblock}

\begin{ansblock}[2章-拓-课时02-T1]
% ans:2章-拓-课时02-T1
\ansitem{17}{(1) θ∈[0,π/2)时斜率由0递增并趋向＋∞；θ∈(π/2,π)时斜率由−∞递增趋向0；(2) 不能，理由见详解．}
\ifansbookdetail
\ansnote{详解}{（1）\(k=\tan\theta\)在\([0,\pi/2)\)上单调递增，\(k\)从\(0\)增大到正无穷大；在\((\pi/2,\pi)\)上也单调递增，\(k\)从负无穷大增大到\(0\)．（2）不能．例如取\(\theta_1=45^{\circ}\)，\(\theta_2=120^{\circ}\)，\(\theta_1<\theta_2\)，但\(k_1=1>k_2=-\sqrt{3}\)，即倾斜角大者斜率反而小；斜率关于倾斜角只在每一段\([0,\pi/2)\)与\((\pi/2,\pi)\)内分别递增，整体上不具有单调性．}
\fi
\end{ansblock}

\begin{ansblock}[2章-拓-课时02-T2]
% ans:2章-拓-课时02-T2
\ansitem{18}{(1) \(k_{AB}=0\)（倾斜角\(0\)）；\(k_{BC}=\sqrt{3}\)（倾斜角\(\pi/3\)）；\(k_{AC}=\sqrt{3}/3\)（倾斜角\(\pi/6\)）；(2) \([\pi/6,\ \pi/3]\)．}
\ifansbookdetail
\ansnote{详解}{（1）\(k_{AB}=(1-1)/(1+1)=0\)，\(k_{BC}=(\sqrt{3}+1-1)/(2-1)=\sqrt{3}\)，\(k_{AC}=(\sqrt{3}+1-1)/(2+1)=\sqrt{3}/3\)．因斜率等于倾斜角的正切值且倾斜角\(\alpha\in[0,\pi)\)，故\(AB\)、\(BC\)、\(AC\)的倾斜角依次为\(0\)、\(\pi/3\)、\(\pi/6\)．（2）当直线\(CD\)绕点\(C\)由\(CA\)逆时针旋转到\(CB\)时，\(CD\)与线段\(AB\)恒有交点（即\(D\)在\(AB\)上），此时\(k_{CD}\)由\(k_{AC}\)增大到\(k_{BC}\)，\(k_{CD}\in[\sqrt{3}/3,\sqrt{3}]\)，对应倾斜角\(\alpha\in[\pi/6,\pi/3]\)．}
\fi
\end{ansblock}

\begin{ansblock}[2章-拓-课时02-T3]
% ans:2章-拓-课时02-T3
\ansitem{19}{(1) 0°；(2) 90°；(3) 45°．}
\ifansbookdetail
\ansnote{详解}{（1）纵坐标同为\(c\)，直线与\(x\)轴平行，倾斜角为\(0^{\circ}\)（\(k=0\)）；（2）横坐标同为\(a\)，直线垂直于\(x\)轴，斜率不存在，倾斜角为\(90^{\circ}\)；（3）\(k=[(c+a)-(b+c)]/(a-b)=(a-b)/(a-b)=1\)（\(a\neq b\)），\(\tan\alpha=1\)，\(\alpha\in[0^{\circ},180^{\circ})\)，故倾斜角为\(45^{\circ}\)．}
\fi
\end{ansblock}

\begin{ansblock}[2章-拓-课时02-T4]
% ans:2章-拓-课时02-T4
\ansitem{20}{[−1/6, 5/3]}
\ifansbookdetail
\ansnote{详解}{\((y+1)/(x+1)=[y-(-1)]/[x-(-1)]\)的几何意义是过\(M(x,y)\)与定点\(N(-1,-1)\)两点的直线的斜率，其中点\(M\)在线段\(y=-2x+8\)（\(x\in[2,5]\)）上运动．线段两端点为\((2,4)\)与\((5,-2)\)．记\(k(x)=(y+1)/(x+1)=(-2x+9)/(x+1)=-2+11/(x+1)\)，在\(x\in[2,5]\)上随\(x\)增大而单调递减，故最大值在\(x=2\)处取得：\(k=(4+1)/(2+1)=5/3\)；最小值在\(x=5\)处取得：\(k=(-2+1)/(5+1)=-1/6\)．所以\((y+1)/(x+1)\)的取值范围为\([-1/6,5/3]\)．}
\fi
\end{ansblock}

\begin{ansblock}[2章-导-课时02-G1]
% ans:2章-导-课时02-G1
\ansitem{21}{D}
\ifansbookdetail
\ansnote{详解}{直线\(l\)的向上方向与\(y\)轴正方向成\(30^{\circ}\)角，\(l\)有两种位置（偏左或偏右）：与\(x\)轴正方向所成的角可能为\(90^{\circ}-30^{\circ}=60^{\circ}\)或\(90^{\circ}+30^{\circ}=120^{\circ}\)，故\(l\)的倾斜角为\(60^{\circ}\)或\(120^{\circ}\)，选D．}
\fi
\end{ansblock}

\begin{ansblock}[2章-导-课时02-G2]
% ans:2章-导-课时02-G2
\ansitem{22}{D}
\ifansbookdetail
\ansnote{详解}{A：当\(\alpha=90^{\circ}\)时，直线的斜率不存在，故A不正确；B：虽然直线的斜率为\(\tan\alpha\)，但只有当\(0^{\circ}\leqslant\alpha<180^{\circ}\)时，\(\alpha\)才是此直线的倾斜角，故B不正确；C：当直线与\(x\)轴平行或重合时，\(\alpha=0^{\circ}\)，\(\sin\alpha=0\)，故C不正确；由倾斜角的定义（任意直线都有倾斜角，范围\(0^{\circ}\leqslant\alpha<180^{\circ}\)）与斜率的定义（\(\alpha\neq 90^{\circ}\)时\(k=\tan\alpha\)）知D正确，选D．}
\fi
\end{ansblock}

\begin{ansblock}[2章-导-课时02-G3]
% ans:2章-导-课时02-G3
\ansitem{23}{D}
\ifansbookdetail
\ansnote{详解}{A：取\(\alpha_1=45^{\circ}\)，\(\alpha_2=135^{\circ}\)，则\(k_1=1\)，\(k_2=-1\)，\(k_1>k_2\)，故A错误；B：取\(\alpha_1=\alpha_2=90^{\circ}\)，则\(k_1\)、\(k_2\)都不存在，故B错误；C：取\(k_1=-1\)，\(k_2=1\)，则\(\alpha_1=135^{\circ}\)，\(\alpha_2=45^{\circ}\)，\(\alpha_1>\alpha_2\)，故C错误；D：斜率相等即\(\tan\alpha_1=\tan\alpha_2\)，而\(\alpha_1\)、\(\alpha_2\in[0^{\circ},180^{\circ})\)且均不为\(90^{\circ}\)，正切值相等则角相等，D正确，选D．}
\fi
\end{ansblock}

\begin{ansblock}[2章-导-课时02-G4]
% ans:2章-导-课时02-G4
\ansitem{24}{(1) 90°；(2) 0°；(3) 45°；(4) 135°．}
\ifansbookdetail
\ansnote{详解}{（1）与\(x\)轴垂直，倾斜角为\(90^{\circ}\)；（2）与\(x\)轴平行或重合，规定倾斜角为\(0^{\circ}\)；（3）角平分线\(y=x\)与\(x\)轴正方向成\(45^{\circ}\)角，倾斜角为\(45^{\circ}\)；（4）角平分线\(y=-x\)向上的方向与\(x\)轴正方向成\(135^{\circ}\)角，倾斜角为\(135^{\circ}\)．}
\fi
\end{ansblock}

\begin{ansblock}[2章-导-课时02-G5]
% ans:2章-导-课时02-G5
\ansitem{25}{10°≤α＜100°}
\ifansbookdetail
\ansnote{详解}{由倾斜角的取值范围\(0^{\circ}\leqslant 2\alpha-20^{\circ}<180^{\circ}\)，各边加\(20^{\circ}\)再除以\(2\)，得\(10^{\circ}\leqslant\alpha<100^{\circ}\)．}
\fi
\end{ansblock}

% ---- 尾块（\tailfill[30mm] 定高参——钉档 §三.3：无参在平衡栏呈「内容尾随框」，贴栏底须定高参；实测无参致末页平衡 Overfull\vbox 12.99pt，定高 30mm 三零） ----
\tailfill[30mm]

\end{multicols}
\end{document}
